\section*{C2: Limbajul IMP. Small-step, Big-step}

\subsection{Big-step}

\subsubsection{Exercițiul 1}

Fie $c$ o instrucțiune și $\sigma, \sigma', \sigma'' \in \Sigma$, cu $(c, \sigma) \Downarrow \sigma'$ și $(c, \sigma) \Downarrow \sigma''$, atunci $\sigma' = \sigma''$.


Inducție după demonstrația lui $(c, \sigma) \Downarrow \sigma'$.

\textbf{Cazul 1}
$c = \mathbf{\mathrm{skip}}$

Avem regula
\[
	(\mathrm{\mathbf{skip}}, \sigma) \Downarrow \sigma
\]
Deci

\[
	\begin{cases}
		(c, \sigma) \Downarrow \sigma' \\
		(c, \sigma) \Downarrow \sigma''
	\end{cases}
	\implies
	\begin{cases}
		\sigma' = \sigma \\
		\sigma'' = \sigma
	\end{cases}
\]

Deci $\sigma' = \sigma''$.

\textbf{Cazul 2:} $c = X := a$.

Ultima regulă din ambele demonstrații este regula atribuirii, deci
\[
	\sigma' = \sigma_{X \mapsto e_\sigma(a)}
	\quad\text{și}\quad
	\sigma'' = \sigma_{X \mapsto e_\sigma(a)}.
\]
Prin urmare, $\sigma' = \sigma''$.

\textbf{Cazul 3:} $c = c_0;c_1$.

Există $\rho',\rho'' \in \Sigma$ astfel încât
\[
	\begin{aligned}
		(c_0,\sigma) & \Downarrow\rho',
		             & (c_1,\rho')       & \Downarrow\sigma',  \\
		(c_0,\sigma) & \Downarrow\rho'',
		             & (c_1,\rho'')      & \Downarrow\sigma''.
	\end{aligned}
\]
Aplicând ipoteza de inducție primei perechi de demonstrații, obținem
$\rho'=\rho''$. Aplicând apoi ipoteza de inducție demonstrației lui
$(c_1,\rho')\Downarrow\sigma'$, obținem $\sigma'=\sigma''$.

\textbf{Cazul 4:} $c=\mathbf{if}\ b\ \mathbf{then}\ c_0\
	\mathbf{else}\ c_1$ și $e_\sigma(b)=1$.

Ultima regulă a primei demonstrații este
\[
	\frac{e_\sigma(b)=1 \qquad (c_0,\sigma)\Downarrow\sigma'}
	{(\mathbf{if}\ b\ \mathbf{then}\ c_0\
		\mathbf{else}\ c_1,\sigma)\Downarrow\sigma'}.
\]
Analizăm acum ultima regulă a demonstrației lui
$(c,\sigma)\Downarrow\sigma''$. Există numai două reguli big-step
pentru o instrucțiune \textbf{if}: prima are premisa $e_\sigma(b)=1$,
iar a doua are premisa $e_\sigma(b)=0$. Cum evaluarea expresiei
booleene $b$ în starea fixată $\sigma$ are valoarea unică $1$, nu poate
fi adevărat simultan că $e_\sigma(b)=0$. Așadar, și a doua demonstrație
folosește prima regulă, având premisa
$(c_0,\sigma)\Downarrow\sigma''$. Aplicând ipoteza de inducție celor
două premise
\[
	(c_0,\sigma)\Downarrow\sigma'
	\quad\text{și}\quad
	(c_0,\sigma)\Downarrow\sigma'',
\]
obținem $\sigma'=\sigma''$.

\textbf{Cazul 5:} $c=\mathbf{if}\ b\ \mathbf{then}\ c_0\
	\mathbf{else}\ c_1$ și $e_\sigma(b)=0$.

Ultima regulă a primei demonstrații este
\[
	\frac{e_\sigma(b)=0 \qquad (c_1,\sigma)\Downarrow\sigma'}
	{(\mathbf{if}\ b\ \mathbf{then}\ c_0\
		\mathbf{else}\ c_1,\sigma)\Downarrow\sigma'}.
\]
Pentru a doua demonstrație, ultima regulă trebuie să fie tot regula
corespunzătoare valorii $0$: expresia $b$, evaluată în aceeași stare
$\sigma$, nu poate avea și valoarea $1$. Prin urmare, a doua
demonstrație are premisa $(c_1,\sigma)\Downarrow\sigma''$. Ipoteza de
inducție aplicată premiselor
\[
	(c_1,\sigma)\Downarrow\sigma'
	\quad\text{și}\quad
	(c_1,\sigma)\Downarrow\sigma''
\]
dă $\sigma'=\sigma''$.

\textbf{Cazul 6:} $c=\mathbf{while}\ b\ \mathbf{do}\ c_0$ și
$e_\sigma(b)=0$.

Regula pentru continuarea buclei cere premisa $e_\sigma(b)=1$, dar aici avem
$e_\sigma(b)=0$. Deoarece evaluarea lui $b$ în starea $\sigma$ are o
singură valoare, regula de continuare nu este aplicabilă. Rămâne numai
regula de oprire:
\[
	\frac{e_\sigma(b)=0}
	{(\mathbf{while}\ b\ \mathbf{do}\ c_0,\sigma)\Downarrow\sigma}.
\]
Concluzia acestei reguli arată că starea finală este chiar starea
inițială $\sigma$. Așadar, din prima derivare
$(c,\sigma)\Downarrow\sigma'$ rezultă $\sigma'=\sigma$.

Același argument se aplică celei de-a doua derivări
$(c,\sigma)\Downarrow\sigma''$, deoarece instrucțiunea, starea inițială
și valoarea lui $e_\sigma(b)$ sunt aceleași. Rezultă
$\sigma''=\sigma$. Prin tranzitivitatea egalității,
\[
	\sigma'=\sigma=\sigma'',
\]
deci $\sigma'=\sigma''$.

\textbf{Cazul 7:} $c=\mathbf{while}\ b\ \mathbf{do}\ c_0$ și
$e_\sigma(b)=1$.

Notăm $w:=\mathbf{while}\ b\ \mathbf{do}\ c_0$. Există
$\rho',\rho''\in\Sigma$ astfel încât
\[
	(c_0,\sigma)\Downarrow\rho',
	\quad (w,\rho')\Downarrow\sigma',
	\qquad
	(c_0,\sigma)\Downarrow\rho'',
	\quad (w,\rho'')\Downarrow\sigma''.
\]
Ipoteza de inducție aplicată demonstrațiilor pentru corpul buclei dă
$\rho'=\rho''$. După această identificare, ipoteza de inducție aplicată
demonstrațiilor buclei dă $\sigma'=\sigma''$. \qed

\subsection{Small-step}

\subsubsection{Exercițiul 2}

\textbf{Cum definim formal închiderea reflexiv-tranzitivă
	$\to^*$ a relației $\to$?}

Fie mulțimea configurațiilor
\[
	\mathcal C:=\mathrm{Instr}\times\Sigma.
\]
Relația small-step este $\to\ \subseteq\mathcal C\times\mathcal C$.
Definim mai întâi puterile relației:
\[
	\to^0:=\Delta_{\mathcal C}
	=\{(\gamma,\gamma)\mid\gamma\in\mathcal C\},
\]
iar, pentru orice $n\in\mathbb N$,
\[
	\gamma\to^{n+1}\gamma'
	\iff
	\exists\delta\in\mathcal C:
	\gamma\to^n\delta\ \text{și}\ \delta\to\gamma'.
\]
Atunci definim
\[
	\boxed{\ \to^*:=\bigcup_{n\in\mathbb N}\to^n\ }.
\]
Prin urmare,
\[
	\gamma\to^*\gamma'
	\iff
	\exists n\in\mathbb N,\ \exists\gamma_0,\ldots,\gamma_n\in\mathcal C
\]
astfel încât
\[
	\gamma_0=\gamma,\qquad
	\gamma_n=\gamma',\qquad
	\gamma_i\to\gamma_{i+1}
	\quad(0\leq i<n).
\]

Pentru $n=0$ nu se execută niciun pas, iar condiția devine
$\gamma=\gamma'$. Acesta este cazul reflexiv:
\[
	\gamma\to^*\gamma.
\]
Pentru $n\geq1$, relația descrie o succesiune finită de pași
small-step.

Relația $\to^*$ este tranzitivă: dacă $\gamma\to^*\delta$ prin $m$ pași
și $\delta\to^*\gamma'$ prin $n$ pași, concatenarea celor două șiruri
dă $\gamma\to^*\gamma'$ prin $m+n$ pași.

În plus, $\to^*$ conține relația $\to$, deoarece
$\to=\to^1\subseteq\to^*$. Ea este chiar cea mai mică relație
reflexivă și tranzitivă ce conține $\to$: pentru orice relație
$R\subseteq\mathcal C\times\mathcal C$,
\[
	\Delta_{\mathcal C}\subseteq R,\quad
	{\to}\subseteq R,\quad
	R\circ R\subseteq R
	\quad\Longrightarrow\quad
	{\to^*}\subseteq R.
\]

\subsubsection{Exercițiul 3}

Este continuarea demonstrației de pe penultimul slide din cursul 2.

\subsubsection{Exercițiul 3}

Este continuarea demonstrației de pe penultimul slide din cursul 2.

Tratăm cazul instrucțiunii \textbf{IF}. Presupunem că avem:
\[
	(\instr{if}\ b\ \instr{then}\ c_0\ \instr{else}\ c_{1}, \sigma) \to^* (\instr{skip}, \sigma')
\]

Analizăm cele două cazuri posibile pentru evaluarea condiției booleene $b$ în starea $\sigma$:

\textbf{Cazul 1:} $e_{\sigma}(b) = 1$
Primul pas de reducere conform semanticii small-step este:
\[
	(\instr{if}\ b\ \instr{then}\ c_0\ \instr{else}\ c_{1}, \sigma) \to (c_0, \sigma)
\]
Deoarece știm că execuția ajunge la $\instr{skip}$ în starea $\sigma'$, înseamnă că pașii rămași sunt:
\[
	(c_0, \sigma) \to^* (\instr{skip}, \sigma')
\]
Aplicând \textit{ipoteza inductivă} (deoarece $c_0$ este mai mică structural decât instrucțiunea \textbf{if} inițială), rezultă că:
\[
	(c_0, \sigma) \Downarrow \sigma'
\]
Conform regulilor semanticii big-step, știind că $e_{\sigma}(b) = 1$ și $(c_0, \sigma) \Downarrow \sigma'$, deducem concluzia dorită:
\[
	(\instr{if}\ b\ \instr{then}\ c_0\ \instr{else}\ c_{1}, \sigma) \Downarrow \sigma'
\]

\textbf{Cazul 2:} $e_{\sigma}(b) = 0$
În mod similar, primul pas de reducere small-step este evaluarea ramurii false (atenție, starea rămâne $\sigma$):
\[
	(\instr{if}\ b\ \instr{then}\ c_0\ \instr{else}\ c_{1}, \sigma) \to (c_1, \sigma)
\]
Rezultă că secvența de reduceri continuă astfel:
\[
	(c_1, \sigma) \to^* (\instr{skip}, \sigma')
\]
Conform \textit{ipotezei inductive} aplicate pe sub-instrucțiunea $c_1$, obținem:
\[
	(c_1, \sigma) \Downarrow \sigma'
\]
Folosind regula big-step pentru ramura falsă (unde $e_{\sigma}(b) = 0$ și $(c_1, \sigma) \Downarrow \sigma'$), concluzionăm:
\[
	(\instr{if}\ b\ \instr{then}\ c_0\ \instr{else}\ c_{1}, \sigma) \Downarrow \sigma'
\]

